\section{Discussion}
\label{sec:discussion}

\subsection{What the Completeness Score Measures---and Why 74.0\% Is Informative}

The mean completeness score of 74.0/100 should not be read as a grade for the extracted
models. It should be read as the answer to a question that has never been asked
systematically before: \emph{if you try to reconstruct the compartmental model from a
peer-reviewed paper, how far can you get?}

Prior to this framework, that question had no instrument. Researchers either trusted
that published models were fully specified (they are not) or informally noted missing
details without measuring the gap. A score of 74.0\% means that, on average, the
reconstructed model recovers roughly three-quarters of the gold-standard structure,
traces most of its fills to source text, achieves high numerical parameter accuracy where
parameters can be found, and resolves about half its internal structural errors
automatically. The 26\% shortfall is not a failure of the pipeline---it is a measurement
of what the papers themselves do not specify.

This framing is the primary scientific contribution of the score. A perfect score of 100
would not mean the model is correct; it would mean the paper is fully explicit. A low
score, as in Ebola (55.4), does not mean the Ebola model is wrong; it means the Ebola
paper leaves more implicit. The score is a property of the \emph{paper}, not the
\emph{disease}.

\subsection{Failure Case: Ebola---When Parameters Are Entirely Implicit}

Ebola is the lowest-scoring disease (55.4/100), driven by two exceptional failures:
low traceability ($T = 37.5\%$) and partial gap reduction ($G = 33.3\%$). Most
Phase~3 fills for Ebola were flagged---parameter values were only partially recoverable from
the paper text via RAG or LLM inference.

Examining the paper~\cite{Abbas2024Ebola} reveals why. The Ebola model uses a
conformable fractional-order ODE formulation. Parameter values are stated in terms of
posterior estimates from Bayesian fitting to surveillance data, reported in a
supplementary table with uncertainty intervals. The model text gives ODE structure but
describes parameters only symbolically (``$\alpha_1$ is the progression rate from E to I'')
without attaching a numeric value to any name. The structural validator correctly detected
\emph{parameter layer contamination} in the extracted model because the LLM injected
posterior means into parameter expressions, conflating the estimation result with the
model input.

This reveals a specific pathology in the Ebola paper's writing style: it separates
model \emph{structure} (written in full) from model \emph{calibration} (relegated to
supplementary material in a format not connected by name to the ODE parameters).
This is not unusual for Bayesian epidemiological papers, and it represents exactly
the kind of incompleteness our framework is designed to detect. A reader attempting
to reproduce the Ebola model would face the same obstacle: the paper tells you what
the compartments mean but not what to put in them for a default run.

\subsection{Failure Case: Malaria---When Parameter Notation Resists Matching}

Malaria achieves perfect compartment recall (1.0) and perfect flow recall (1.0) but
parameter recall plateaus at 0.737---even after Phase~3. The 12 remaining gaps after
filling are all parameters.

The malaria paper~\cite{Akowe2025Malaria} integrates vector and non-vector transmission
pathways, introducing composite rate expressions such as
$\Lambda_h = \frac{b \beta_{mh} I_m}{N_h} + \Lambda_{nh}$, where $\Lambda_{nh}$ is
a non-vector force of infection term defined recursively elsewhere. When the RAG index
retrieves such a passage, the fill value is the symbolic expression, not a number. The
Phase~3 completeness score correctly penalizes this: the fill is traced to source
(traceability is 66.7\%) but the numeric accuracy comparison fails because the gold
value is a constant and the fill is a formula.

This failure reveals a general tension in mathematical biology writing: rate parameters
are frequently \emph{defined compositionally}---as functions of other parameters---
rather than as independent constants. A compartmental model extracted from such a paper
necessarily inherits this incompleteness unless the reader also interprets the full
mathematical derivation. This is an inherent limit of extraction from text, and
our framework surfaces it explicitly rather than hiding it behind a high F1 score.

\subsection{Other Challenging Cases: Flu, Tuberculosis, and HIV}

Three other diseases score below 70/100 for distinct but related reasons.

\paragraph{Influenza (63.1/100) --- parameters buried in figures.}
The flu paper~\cite{Corbella2017MonitoringAP} achieves perfect compartment and flow
recall before Phase~3, leaving only one unfilled gap: a transmission parameter whose
numeric value appears only in a figure caption rather than in the main text.
RAG indexes the paper body but cannot extract values from embedded figure graphics,
so the gap reduction score is 0\% despite strong overall structure.
The failure mode here is not a problem with our pipeline---it is a reporting choice in
the paper. A numeric value placed only in a figure is inaccessible to any text-based
extraction system. Future work on figure-level OCR or caption parsing could close
this class of gaps.

\paragraph{Tuberculosis (67.1/100) --- composite ODE coefficients.}
The TB paper~\cite{tuberculosis} defines its transmission parameters as composed
expressions embedded in the ODE system (e.g., $\beta_{t} = \beta_0 (1 - \varepsilon)$),
where each component is defined in a different subsection.
RAG retrieves individual passages but cannot assemble multi-part definitions spread
across sections, so gap reduction reaches only 25\%.
Fill traceability is high (100\%) because the fills that \emph{are} made are correctly
sourced; the missing gaps are simply ones where no single passage contains a complete value.
This reflects a writing practice common in TB and multi-strain models: parameters are
introduced incrementally rather than in a consolidated table.

\paragraph{HIV (81.6/100) --- structural complexity outpaces repair.}
HIV scores well on recall (81.6 after Phase~3) but has the lowest structural integrity
(25.0\%) of any disease. The HIV model~\cite{Espitia2022} is a three-risk-group formulation
with 14+ group-specific parameters ($b_{hw}$, $b_{wm}$, $c_{hm}$, \ldots).
After Phase~3 fills compartments and flows correctly, the structural repairer cannot
automatically assign each group-specific rate to its correct flow without a formal
parameter-to-flow mapping table that the paper does not provide.
The 25\% structural integrity score is therefore not a gap in extraction --- the
entities are found --- but a gap in the paper's formal specification of how those
entities connect.

\paragraph{Common root cause.}
These cases share a pattern: the paper's \emph{informal} text contains enough information
to understand the model if read by a domain expert, but that information is not organized
in a way that supports automated extraction. Values in figures, multi-part definitions,
and implicit parameter assignments all require reader inference that goes beyond what any
text retrieval system can provide today.

\subsection{What Structural Errors Reveal About How Modelers Write Papers}

The 192 structural errors detected across 10 Phase~3 models are not random noise.
Each error class corresponds to a recurring pattern in how compartmental modeling
papers are written:

\paragraph{Uniform parameter collapse (10 occurrences, 100\% resolved).}
This error---where a single parameter is assigned to all flows---occurs when the paper
introduces rate parameters symbolically ($\beta$, $\gamma$, $\mu$) but discusses each
one in isolation without explicitly stating which transition each governs. The LLM,
lacking an explicit assignment table, defaults to reusing the first extracted parameter
everywhere. Papers that provide a single ``parameter table'' without a flow-by-flow
breakdown systematically trigger this error. The 100\% resolution rate confirms that
the error is deterministic given a correct parameter semantic classification.

\paragraph{Orphaned parameters (88 occurrences, 32\% resolved).}
This is the most common error class and the hardest to resolve. An orphaned parameter
appears in the model's parameter list but is not connected to any flow. It arises
because papers often declare a large parameter set in a table and then use only a
subset in the differential equations, with the remaining parameters entering only in
derived quantities (e.g., the basic reproduction number $R_0$). When the LLM faithfully
copies all declared parameters into the model but does not have enough context to map
each one to a specific flow, orphaned parameters result. This is a direct consequence
of the gap between \emph{parameter inventory} (what the paper lists) and \emph{parameter
deployment} (what the paper's ODEs actually use). HIV (14 orphaned params) and Malaria
(11 orphaned) are the most affected because their multi-group and multi-pathway structures
declare group-specific parameters ($b_{hw}$, $c_{hm}$, $\Lambda_{nh}$) whose assignment
to specific flows is only recoverable from the ODE system itself, not the text.

\paragraph{Self-referential flows (12 occurrences, 25\% resolved).}
A self-referential flow has the same compartment as source, target, and contact
compartment---an impossible mass-action term. This arises when a paper writes
``susceptibles are infected at rate $\beta S I / N$'' without explicitly naming the target
state. The LLM correctly identifies the flow as a contact process but resolves the target
to the same compartment because no Exposed or Infectious compartment is mentioned in
that passage. Papers that introduce S, I, R sequentially but defer the transition
diagram to a figure (rather than stating ``S transitions to I'') trigger this pattern.
The 25\% resolution rate reflects that only models with an explicit intermediate
compartment can be structurally repaired; simple SIR models with no latent stage must
remain annotated.

\paragraph{Parameter layer contamination (24 occurrences, 33\% resolved).}
This error occurs when the LLM inserts posterior estimates, likelihood values, or
confidence intervals into parameter expressions---conflating Bayesian model fitting
output with ODE model inputs. It is entirely a consequence of papers that present
both the model and its fitted parameters in the same section, using identical notation
for prior and posterior quantities. The relatively low resolution rate (33\%) reflects
that clearing a contaminated expression removes information the paper provides but
does not supply the correct deterministic value.

\paragraph{Aggregate insight.}
Taken together, these four patterns point to the same root cause: the convention in
epidemiological papers of separating the structural description (compartment diagram,
ODEs) from the operational description (parameter values, assignments, initial conditions)
without explicitly linking them. This convention is fine for readers who can do the
linking mentally, but it is precisely what makes automated model extraction hard and
makes formal reproducibility assessment necessary. Our framework is designed to detect
and quantify this gap, not to judge the papers that exhibit it.

\subsection{RAG vs.\ LLM Inference: Two Roles for the Same Goal}

When evaluated by \emph{recall}---the metric that measures coverage without penalizing
extra entities---LLM-only is nearly as effective as RAG-only (mean $\bar{R}$ of 0.936
vs.\ 0.951, Table~\ref{tab:ablation}). Both modes substantially outperform the
Phase~2 draft (0.846). The key difference is what each mode adds when it succeeds:

RAG fills are \emph{traceable}---they are sourced from a specific passage in the paper's
own text, which is exactly the kind of evidence a completeness audit should require.
When RAG fills a missing flow, we can point to the sentence that specifies it.
LLM-inferred fills are \emph{plausible but unverifiable}---the LLM generates a
structurally reasonable entity based on epidemiological conventions, but that entity
may not correspond to anything the authors intended. LLM fills are flagged with
\texttt{source=inference} and lower the fill traceability component of the completeness
score accordingly.

The practical implication has two sides. As a \emph{coverage tool} (``what entities
does this model need?''), LLM inference adds real value: it reaches nearly the same
recall as RAG and costs no additional paper reading. As an \emph{audit tool} (``what
did the paper specify?''), only RAG-sourced fills are trustworthy evidence of the paper's
intent. A gap that RAG cannot fill but LLM can represents a genuine specification
failure in the paper---the model element exists in the final model but was never written
down explicitly. That is precisely the kind of failure our framework is designed to
surface and quantify.
